% Figure 3: where the money goes. Token classes are what Eq. (9) prices
% differently, so the figure is drawn in those classes rather than in a single
% token total. Coordinates from figures/data/token_classes.dat.

\begin{figure}[t]
\centering
\begin{tikzpicture}[font=\small]
\begin{axis}[
  width=0.92\linewidth, height=5.4cm,
  ybar, bar width=7pt,
  ymode=log, log origin=infty,
  ymin=1, ymax=2200,
  ytick={1,10,100,1000},
  yticklabels={1,10,100,1000},
  ylabel={токены (миллионы, лог. шкала)},
  xtick={0,1,2,3},
  xticklabels={контроль,$k{=}0$,$k{=}2$,$k{=}6$},
  enlarge x limits=0.16,
  grid=major, grid style={black!12},
  legend style={at={(0.98,0.97)}, anchor=north east, font=\scriptsize,
                draw=black!25, fill=white, fill opacity=0.92, text opacity=1},
  legend cell align=left,
]
\addplot[draw=black!60, fill=green!25] table[x=x, y=cached]
  {figures/data/token_classes.dat};
\addlegendentry{кэшированный вход (дешевле всего, $\gamma$)}
\addplot[draw=black!60, fill=blue!30] table[x=x, y=uncached]
  {figures/data/token_classes.dat};
\addlegendentry{некэшированный вход ($\alpha$)}
\addplot[draw=black!60, fill=red!45] table[x=x, y=output]
  {figures/data/token_classes.dat};
\addlegendentry{выход (дороже всего, $\beta$)}
\end{axis}
\end{tikzpicture}
\caption{Потребление токенов по тарифицируемым классам по всей серии,
декомпозиция, которую тарифицирует Ур.~\eqref{eq:cost}. Столбцы
сгруппированы, а не составлены друг на друга, а ось логарифмическая, потому
что классы различаются на два порядка. Контрольная группа потребляет
абсолютно больше всего токенов, но \TotCacheBase{}М из них — кэшированные
чтения, тарифицируемые по самой дешёвой ставке: только дополняемый промпт
почти полностью обслуживается из кэша префикса. Руки состояния потребляют
намного меньше токенов и платят за них больше, потому что блок состояния,
переписываемый на каждом ходу, теряет этот кэш (вплоть до \TotCacheTailSix{}М)
и потому что поддержание состояния переносит токены в класс выхода. Именно
поэтому сравнение по числу токенов и сравнение по стоимости расходятся.}
\label{fig:tokenclasses}
\end{figure}
